% --- Overleaf / paste-in notes (delete this block after reading) ---
% - Place after \appendix in your main file, or merge into a single supplement.
% - Preamble: \usepackage{amsmath,booktabs} (booktabs only if you add tables here).
% - This file does not use \input or \ref to your repo; adjust the section title
%   or \label{app:oklahoma} if it clashes with your document.
% ---------------------------------------------------------------------------

\section{Oklahoma application: supplementary detail}\label{app:oklahoma}

\subsection{Relation to the computational supplement}

The analysis is documented in full in a separate computational report (Quarto/HTML) generated from the same fitted results archive as the paper. That report defines an all-or-nothing estimand $\Delta_{\mathrm{AoN}}$ on the county partition: the difference in total expected event counts under universal treatment versus universal control. The application section of the paper uses the notation $\hat\Delta_{\mathrm{AoI}}$ for the same scientific contrast over the post-directive window. The numerical results reported for the naive versus SEM comparison correspond to the bivariate ETAS models with nonparametric (KDE) background and all triggering parameters estimated (models E and F in the supplement), evaluated at the horizon given by the stored analysis configuration.

\subsection{Study design}

\begin{itemize}
  \item \textbf{Treatment.} Oklahoma Corporation Commission directive AOI\_20150318, requiring reduced wastewater disposal within the designated Area of Interest.
  \item \textbf{Spatial units.} US Census counties in Oklahoma; counties whose centroid falls inside the AOI are assigned to treatment, remainder to control. The supplement lists treated counties, tile counts, and event counts by phase.
  \item \textbf{Time windows.} Pre-treatment and post-treatment intervals are anchored at the directive time $t^*$ and at the catalog end date recorded in the study metadata.
  \item \textbf{Magnitude threshold.} Only events with magnitude at least $m_0$ are used, with $m_0$ fixed by the analysis configuration (consistent with the catalog design file).
  \item \textbf{Pre-treatment split for the background.} The first half of pre-treatment events is reserved to train the spatial background (KDE on control, pre-$t^*$); the second half is used together with post-$t^*$ data in estimation with temporal carry-over, as described in the supplement under the nonparametric background discussion.
\end{itemize}

\subsection{Model families}

Besides the naive (E) and SEM (F) bivariate ETAS fits with KDE background and freely estimated triggering parameters, the supplement fits additional specifications for comparison: homogeneous backgrounds, structurally constrained parameters, alternative KDE variants, and univariate ETAS$+$KDE models. It also refits the E/F class under alternative spatial partitions (grids scaled to an estimated triggering range, and a coarse partition with the AOI as a single treated region) and under a small set of KDE bandwidth multipliers.

\subsection{Bivariate ETAS intensity}

Let $z \in \{0,1\}$ denote latent process membership (control or treated). The conditional intensity for process $k \in \{0,1\}$ is
\begin{equation*}
\lambda_k(t,x,y)
  = \frac{\mu_k}{|S_k|}\, W_k(x,y)
  + \sum_{\ell \in \{0,1\}} \sum_{\substack{j : t_j < t \\ z_j = \ell}}
    A_{k\ell}\, e^{\alpha_{m,k\ell}(m_j - m_0)}\,
    g(t - t_j)\, f(x-x_j, y-y_j \mid m_j),
\end{equation*}
where $A_{01}$ and $A_{10}$ capture cross-excitation. For models E and F, $W_k$ is a normalized kernel density estimate from held-out control pre-$t^*$ events; for homogeneous-background models in the supplement, $W_k \equiv 1$. The temporal kernel $g$ (Omori--Utsu) and spatial kernel $f$ (power law in distance with magnitude-dependent scale) are the standard ETAS forms used in the implementation and are written explicitly in the supplement.

\subsection{Temporal truncation}

Likelihood evaluation and simulation omit triggering contributions from parents more than $t_{\mathrm{trunc}}$ days before each event. The supplement tabulates $t_{\mathrm{trunc}}$, structural parameters $(c,p,q)$, and the implied fraction of Omori tail mass retained within the truncation window. The same rule is applied in estimation and in counterfactual simulation.

\subsection{Stochastic EM}

The SEM alternates stochastic updates to latent event labels with re-estimation of parameters. The supplement reports outer and inner iteration counts, numbers of proposals and retained labellings, selection temperature, and change-factor bounds, together with optional pilot tuning for the control-fixed SEM variant. Diagnostic figures include traces of the adaptive objective, cumulative label flips, and confusion matrices comparing location-based labels to SEM-implied labels for each SEM fit.

\subsection{Counterfactual means and bootstrap}

Expected counts under ``control everywhere'' and ``actual regime,'' and their difference, are obtained by Monte Carlo over the stored all-or-nothing simulation draws. The supplement gives a full table of these quantities (and simulation dispersion) for every fitted model, and, when bootstrap output is available, bootstrap standard deviations and recentred quantile-type intervals for the primary E/F pair.

The bootstrap histogram in the main text uses bias-adjusted replicate means: replicate means are shifted so that, for each model, the mean of the bootstrap cloud matches the model-implied mean total saved under the actual regime before comparison across models. The supplement additionally shows empirical CDFs of bootstrap draws, running means over replicate index, and histograms of Monte Carlo draws of total saved. Those extra figures can be exported as PDFs alongside the main-text figure if needed for the journal.

\subsection{Content available in the supplement but not in the main text}

\begin{description}
  \item[Exploratory figures.] Magnitude histogram for events above $m_0$; cumulative event time series (equivalent to the cumulative-count figure in the paper); map of county treatment assignment; side-by-side pre- and post-$t^*$ spatial event maps on the county tessellation.
  \item[Background estimation.] Plot of the fitted KDE background intensity (when saved) and a table of KDE training sample sizes, bandwidth, and estimated triggering range.
  \item[Structural parameters.] Tables listing fixed versus estimated ETAS structural parameters and, when available, control-only snapshot fits at reference times.
  \item[All models.] One table collating observed post-$t^*$ counts, simulated means and standard deviations under each regime, expectations under ``treated everywhere,'' and bootstrap summaries for the primary pair.
  \item[E/F parameters.] Full vector of estimated ETAS parameters for models E and F side by side.
  \item[Simulated ATE distributions.] Faceted histograms of Monte Carlo total saved across model families, and analogous figures stratified by spatial partition.
  \item[SEM diagnostics.] Time series of the adaptive score, cumulative relabelling activity, and confusion matrices for SEM fits.
  \item[Partition and bandwidth sensitivity.] Maps of alternative partitions; a large table comparing expected counts, saved events, spillover summaries, and marginal parameters across partitions; line plots of marginal parameters versus partition; a table of results across KDE bandwidth multipliers on the county partition.
  \item[Reproducibility.] A final listing of configuration scalars used in the run.
\end{description}
